\begin{circuitikz}[
    x=1cm,
    y=1cm,
    every node/.style={font=\footnotesize\sffamily},
    net/.style={BookBlue,very thick},
    supply/.style={BookRed,very thick},
    passive/.style={
      draw=BookAmber,
      fill=BookAmber!15,
      rounded corners=0.6mm,
      minimum width=1.55cm,
      minimum height=0.55cm,
      align=center
    }
  ]
  \node[font=\bfseries\Large\sffamily\color{BookNavy}] at (10.7,10.2)
    {I2C EEPROM：两根总线、两个上拉和八个封装脚};

  \draw[BookNavy,very thick,fill=BookNavy!9,rounded corners=1.5mm]
    (0.45,5.45) rectangle (5.55,9.45);
  \node[font=\bfseries\large\sffamily,align=center] at (3.0,8.85)
    {Intel N-Series SoC\\I2C0};
  \draw[BookNavy,thick] (4.95,7.65) -- (5.55,7.65)
    node[ocirc,pos=1,fill=white]{};
  \node[anchor=east,align=right] at (4.80,7.65)
    {\textbf{ball CW42}\\\code{I2C0\_SCL}};
  \draw[BookNavy,thick] (4.95,6.45) -- (5.55,6.45)
    node[ocirc,pos=1,fill=white]{};
  \node[anchor=east,align=right] at (4.80,6.45)
    {\textbf{ball CU42}\\\code{I2C0\_SDA}};

  \draw[BookTeal,very thick,fill=BookTeal!8,rounded corners=1.5mm]
    (13.0,5.10) rectangle (18.85,9.65);
  \node[font=\bfseries\large\sffamily,align=center] at (15.93,9.05)
    {U4 24C02\\I2C EEPROM};
  \draw[BookTeal,thick] (13.0,7.65) -- (13.6,7.65)
    node[ocirc,pos=0,fill=white]{};
  \node[anchor=west] at (13.75,7.65) {\textbf{pin 6}\quad\code{SCL}};
  \draw[BookTeal,thick] (13.0,6.45) -- (13.6,6.45)
    node[ocirc,pos=0,fill=white]{};
  \node[anchor=west] at (13.75,6.45) {\textbf{pin 5}\quad\code{SDA}};
  \draw[net,{Latex[length=2mm]}-{Latex[length=2mm]}]
    (5.55,7.65) -- node[above,font=\scriptsize]{\code{I2C0\_SCL}} (13.0,7.65);
  \draw[net,{Latex[length=2mm]}-{Latex[length=2mm]}]
    (5.55,6.45) -- node[above,font=\scriptsize]{\code{I2C0\_SDA}} (13.0,6.45);

  \node[passive] (rp1) at (7.70,8.75) {R3\\\(\SI{2.2}{k\ohm}\)};
  \node[passive] (rp2) at (10.80,8.75) {R4\\\(\SI{2.2}{k\ohm}\)};
  \draw[net] (7.70,7.65) -- (rp1.south);
  \draw[net] (10.80,6.45) -- (rp2.south);
  \draw[supply] (rp1.north) -- (7.70,9.45) -- (10.80,9.45) -- (rp2.north);
  \node[font=\bfseries\color{BookRed},above] at (9.25,9.45) {\code{VCC\_IO}};

  \draw[BookTeal,thick] (18.85,8.35) -- (19.40,8.35)
    node[ocirc,pos=0,fill=white]{};
  \node[anchor=east] at (18.70,8.35) {\textbf{pin 8 VCC}};
  \draw[supply] (19.40,8.35) -- (20.30,8.35)
    node[right,font=\bfseries]{\code{VCC\_IO}};

  \draw[BookTeal,thick] (18.85,7.25) -- (19.40,7.25)
    node[ocirc,pos=0,fill=white]{};
  \node[anchor=east] at (18.70,7.25) {\textbf{pin 7 WP}};
  \draw[BookNavy,thick] (19.40,7.25) -- (20.00,7.25) node[ground]{};

  \draw[BookTeal,thick] (18.85,6.15) -- (19.40,6.15)
    node[ocirc,pos=0,fill=white]{};
  \node[anchor=east,align=right] at (18.70,6.15)
    {\textbf{pins 1/2/3}\\\code{A0/A1/A2}};
  \draw[BookNavy,thick] (19.40,6.15) -- (20.00,6.15) node[ground]{};

  \draw[BookTeal,thick] (18.85,5.35) -- (19.40,5.35)
    node[ocirc,pos=0,fill=white]{};
  \node[anchor=east] at (18.70,5.35) {\textbf{pin 4 GND}};
  \draw[BookNavy,thick] (19.40,5.35) -- (20.00,5.35) node[ground]{};

  \node[font=\bfseries\Large\sffamily\color{BookNavy}] at (10.7,4.45)
    {UART 调试口：发送脚接对方接收脚};
  \draw[BookNavy,very thick,fill=BookNavy!9,rounded corners=1.5mm]
    (0.45,0.10) rectangle (7.20,3.85);
  \node[font=\bfseries\large\sffamily,align=center] at (3.83,3.25)
    {Intel N-Series SoC\\UART0 pad mux};
  \draw[BookNavy,thick] (6.60,2.30) -- (7.20,2.30)
    node[ocirc,pos=1,fill=white]{};
  \node[anchor=east,align=right] at (6.45,2.30)
    {\textbf{ball DB40}\\\code{GPP\_H11/UART0\_TXD}};
  \draw[BookNavy,thick] (6.60,1.10) -- (7.20,1.10)
    node[ocirc,pos=1,fill=white]{};
  \node[anchor=east,align=right] at (6.45,1.10)
    {\textbf{ball DB34}\\\code{GPP\_H10/UART0\_RXD}};

  \draw[BookAmber!80!black,very thick,fill=BookAmber!12,rounded corners=1.5mm]
    (15.0,-0.15) rectangle (21.0,4.05);
  \node[font=\bfseries\large\sffamily,align=center] at (18.0,3.45)
    {J2 调试排针\\接电平匹配 USB--UART};
  \draw[BookAmber!80!black,thick] (15.0,2.30) -- (15.60,2.30)
    node[ocirc,pos=0,fill=white]{};
  \node[anchor=west] at (15.75,2.30)
    {\textbf{pin 2}\quad\code{RXD(adapter)}};
  \draw[BookAmber!80!black,thick] (15.0,1.10) -- (15.60,1.10)
    node[ocirc,pos=0,fill=white]{};
  \node[anchor=west] at (15.75,1.10)
    {\textbf{pin 3}\quad\code{TXD(adapter)}};
  \draw[net,-{Latex[length=2mm]}] (7.20,2.30) --
    node[above,font=\scriptsize]{SoC TX \(\rightarrow\) adapter RX} (15.0,2.30);
  \draw[net,{Latex[length=2mm]}-] (7.20,1.10) --
    node[above,font=\scriptsize]{SoC RX \(\leftarrow\) adapter TX} (15.0,1.10);

  \draw[BookAmber!80!black,thick] (18.0,-0.15) -- (18.0,-0.55)
    node[ocirc,pos=0,fill=white]{};
  \node[anchor=west] at (18.15,-0.35) {\textbf{pin 1 GND}};
  \draw[BookNavy,thick] (18.0,-0.55) node[ground]{};

  \node[
    draw=BookMuted,
    fill=white,
    rounded corners=1mm,
    align=left,
    text width=5.1cm,
    font=\scriptsize\sffamily
  ] at (10.90,0.05) {
    先确认 UART pad 的实际 I/O 电压，再选择匹配的 USB--UART 适配器。
    这里画的是 CMOS/TTL 串口，不得直接连接使用正负电压的 RS-232 DB9。
  };
\end{circuitikz}
